% 51e Olympiade Internationale de Mathématiques — Astana, Kazakhstan.
% Énoncés officiels en français (imo-official.org). Chaque problème vaut 7 points.

\begin{center}
  \emph{Premier jour --- mercredi 7 juillet 2010 --- Durée : 4 heures 30}
\end{center}

\subsection*{Problème 1}
Déterminer toutes les fonctions $f : \mathbb{R} \to \mathbb{R}$ telles que, pour tous $x, y \in \mathbb{R}$ :
\[
f\bigl(\lfloor x \rfloor y\bigr) = f(x)\bigl\lfloor f(y) \bigr\rfloor.
\]
(On note $\lfloor z \rfloor$ le plus grand entier inférieur ou égal à $z$.)

\subsection*{Problème 2}
Soit $I$ le centre du cercle inscrit dans le triangle $ABC$ et soit $\Gamma$ son cercle circonscrit. La droite
$(AI)$ recoupe $\Gamma$ en $D$. Soit $E$ un point de l'arc $\overset{\frown}{BDC}$ et $F$ un point du côté
$[BC]$ tels que
\[
\widehat{BAF} = \widehat{CAE} < \tfrac{1}{2}\widehat{BAC}.
\]
Soit enfin $G$ le milieu du segment $[IF]$.

Montrer que les droites $(DG)$ and $(EI)$ se coupent en un point de $\Gamma$.

\subsection*{Problème 3}
$\mathbb{N}^*$ désigne l'ensemble des entiers strictement positifs.

Déterminer toutes les fonctions $g : \mathbb{N}^* \to \mathbb{N}^*$ telles que, pour tous
$m, n \in \mathbb{N}^*$,
\[
\bigl(g(m) + n\bigr)\bigl(m + g(n)\bigr)
\]
soit un carré parfait.

\bigskip
\begin{center}
  \emph{Second jour --- jeudi 8 juillet 2010 --- Durée : 4 heures 30}
\end{center}

\subsection*{Problème 4}
Soit $P$ un point intérieur au triangle $ABC$. Les droites $(AP)$, $(BP)$ et $(CP)$ recoupent $\Gamma$, cercle
circonscrit au triangle $ABC$, respectivement aux points $K$, $L$ et $M$. La tangente en $C$ à $\Gamma$ coupe
la droite $(AB)$ en $S$. On suppose que $SC = SP$.

Montrer que $MK = ML$.

\subsection*{Problème 5}
Au début, chacune des six boîtes $B_1, B_2, B_3, B_4, B_5, B_6$ contient un jeton. Deux types d'opération
sont possibles :
\begin{itemize}
  \item[\emph{Type} 1 :] Choisir une boîte non vide $B_j$ avec $1 \leqslant j \leqslant 5$ ; ôter un jeton de la
  boîte $B_j$ et ajouter deux jetons à la boîte $B_{j+1}$.
  \item[\emph{Type} 2 :] Choisir une boîte non vide $B_k$ avec $1 \leqslant k \leqslant 4$ ; ôter un jeton de la
  boîte $B_k$ et échanger les contenus des boîtes (éventuellement vides) $B_{k+1}$ et $B_{k+2}$.
\end{itemize}
Est-il possible, à la suite d'un nombre fini de telles opérations, que les boîtes $B_1, B_2, B_3, B_4, B_5$
soient vides et que la boîte $B_6$ contienne exactement $2010^{2010^{2010}}$ jetons ? (Noter que
$a^{b^c} = a^{(b^c)}$.)

\subsection*{Problème 6}
Soit $a_1, a_2, a_3, \ldots$ une suite de nombres réels strictement positifs. On suppose qu'il existe un
entier $s$ strictement positif tel que, pour tout $n > s$, on ait :
\[
a_n = \max\{a_k + a_{n-k} \mid 1 \leqslant k \leqslant n - 1\}.
\]
Montrer qu'il existe des entiers strictement positifs $\ell$ et $N$, avec $\ell \leqslant s$, tels que, pour
tout $n \geqslant N$ :
\[
a_n = a_\ell + a_{n-\ell}.
\]
